\section{User Study Methodology}
\label{sec:method}

% \subsection{Study type and participants}

This work is an experience report on the design and qualitative evaluation of a web-based educational game that supports the understanding of Nielsen's ten heuristics. The artifact was analyzed by HCI students, professors, and industry professionals, who were reached via convenience sampling. 

% deployed in two complementary settings. The first was an HCI evaluation course at a Brazilian university (omitted for double-blind review), where the course instructor applied the game during regular class time, after the ten heuristics had already been introduced through the standard course content. This setting closely mirrors the intended classroom use of the artefact, in which the game supports the understanding of heuristics that the instructor has just taught. The second setting was remote participation reached through the research team's network, which broadens the range of participant profiles beyond the host course.

Participants formed a mixed audience of undergraduate and graduate students from computing-related programs, with additional professionals and researchers reached through the wider network. Prior familiarity with Nielsen's heuristics varied among participants: some had been introduced to them this semester, while others encountered them for the first time through the game. This variation is captured by a dedicated item in the demographic questionnaire so that responses can be read against prior exposure when relevant. Recruitment was non-probabilistic; no personally identifiable information was collected. The final sample includes $N = 30$ responses; no entries were removed during the data cleaning (see Section~\ref{sec:results}).

\subsection{Procedure}

Each participant followed the same sequence. They opened the game in a web browser, played end-to-end through the ten heuristics, and then clicked through to a linked questionnaire from the game's final results screen \footnote{The questionnaire was administered through Google Forms in Portuguese, the language approved by the ethics committee for this .}. Its first page presented the informed-consent statement, which had to be accepted before continuing. The remaining pages contained five categorical demographic questions and five open-ended questions about the participant's experience.

The demographic questions covered current role, age, primary area of work or study, prior familiarity with Nielsen's heuristics, and practical experience with HCI or usability evaluation. The open-ended questions were related to (a) overall appreciation of the game, (b) perceived support to understand the heuristics, (c) the effect of gamification and narration on experience and learning, (d) negative points and suggested improvements, and (e) further comments. Responses were automatically collected in a linked spreadsheet available exclusively to the research team.

% The questionnaire deliberately omits Likert-style agreement scales. This choice reflects two design constraints. The expected sample size is small, which limits what numeric scales can reveal at this scale. The research question also targets the mechanisms behind perceived support --- something better captured by open-ended responses that participants are free to phrase in their own terms.

\subsection{Data analysis}

The open-ended responses were analyzed through qualitative content analysis \cite{bardin2011}, following the four-step procedure adopted for this study. First, the responses were screened for engagement: entries that were obviously not engaged with the question (for instance, single-character or filler responses) would be removed at this step. Second, the demographic answers were summarized through charts, each accompanied by a textual description so that the argument can be followed without consulting the figure. Third, the open-ended responses were grouped by question and inductively clustered into themes: similar responses were brought together and named, with cross-participant comparisons to identify shared patterns rather than listing individual answers. Fourth, the themes were exemplified through anonymized excerpts, with participants identified only as $P_1, P_2, \ldots$; where context matters for a given quote, a short demographic tag is added (for example, ``$P_3$ (graduate student, computer science, no prior familiarity with the heuristics)'').

\subsection{Ethical considerations}
\label{sec:ethics}

This experience report respects the Brazilian regulatory framework for research involving human participants. The study is part of a project approved by the host university's Research Ethics Committee, approved under CAAE number (omitted for double-blind review).

Participation was voluntary, non-remunerated, and could be withdrawn at any moment without consequence. Before answering the questionnaire, each participant read an Informed Consent Form, describing the details of the research, and consent had to be explicitly given before any further questions could be answered.
